%% kapa: exam=2 round=22 year=2011 period=1 subject=appraisal_practice subject_ko=감정평가실무 seq=3 number=3 points=20
%% source: Q-Net 감정평가사 기출문제 (cst00309) artlSeq=5040389 file=extracted/1008/5040389/2011년 제22회 감정평가사 2차시험 기출문제/1교시.hwp
%% conversion: kordoc 4.24.51 -> tools/gen2.py (auto); review=proofread-sonnet; flags=none
\begin{longquestion}{【문제 3】}
D건설회사는 총 720세대 규모의 아파트단지 조성사업을 시행하여 입주가 완료되었으나 그 중 12세대는 일반적인 아파트와는 달리 거실 유리창의 일부가 감소되도록 설계되어 입주 후 가치하락액을 산정하여 환불해주기로 하고 환불대상세대 및 환불액 결정을 K감정평가법인에 의뢰하였다. 아래에 제시된 조건과 자료의 범위내에서 K감정평가법인이 수행해야 할 환불대상세대 결정 및 대상세대의 최종 환불액을 평가하시오. (20점) \par
<자료 1> 기본적 사항 \par
1) 환불액 평가의 가격시점은 2011. 08. 01로 한다. \par
2) 환불액 지급대상세대는〈자료 2〉에 제시된 세대 중 연속일조시간이 2시간 미만이고 총일조시간이 4시간 미만인 세대만을 대상으로 한다. \par
3) 환불액은 일조시간을 기초로 산정한 가치하락액을 기준으로 결정한다. 이 경우 본 아파트단지에서 일조권 가치가 전체가치에서 차지하는 비율은 평형에 관계없이 6\%이며 총일조시간(x분)과 해당세대의 가치하락율(y)간의 관계는 다음 산식으로 산정한다. \par
y = 0.06 (1 - x / 240 ) \par
4) 환불대상세대 중 1년 이내에 거래사례가 있는 경우에는 거래사례에 의해 산정한 가치하락액과 일조시간을 기준으로 산정한 가치하락액을 비교하여 적은 금액으로 환불액을 결정한다. \par
<자료 2> 대상 아파트단지 개요 \par
- 소재지 : S시 A구 \par
- 규모 : 총 720세대(10개동 × 각동 72세대) \par
- 층수 : 각동 공히 18개층 높이이며 각층별 세대수는 동일하게 건축되었음 \par
- 동별 현황 \par
ㆍ101동 ～ 108동 : 전세대 85m\textsuperscript{2}형 \par
ㆍ109동, 110동 : 전세대 110m\textsuperscript{2}형 \par
<자료 3> 창면적 감소세대 현황 \par
{\small\centering
\begin{tabular}{|c|c|c|c|c|}
\hline
\makecell{동번호} & \makecell{해당세대} & \makecell{창면적감소비율\newline{}(\%)} & \makecell{총 일조시간\newline{}(분)} & \makecell{연속일조시간\newline{}(분)} \\
\hline
\makecell{101} & \makecell{301호\newline{}302호\newline{}401호\newline{}402호} & \makecell{45\newline{}18\newline{}45\newline{}18} & \makecell{165\newline{}265\newline{}170\newline{}270} & \makecell{95\newline{}183\newline{}98\newline{}185} \\
\hline
\makecell{102} & \makecell{602호\newline{}702호} & \makecell{60\newline{}25} & \makecell{160\newline{}250} & \makecell{93\newline{}170} \\
\hline
\makecell{109} & \makecell{301호\newline{}302호\newline{}401호\newline{}402호} & \makecell{45\newline{}18\newline{}45\newline{}18} & \makecell{165\newline{}265\newline{}170\newline{}270} & \makecell{125\newline{}183\newline{}128\newline{}185} \\
\hline
\makecell{110} & \makecell{602호\newline{}702호} & \makecell{60\newline{}25} & \makecell{160\newline{}250} & \makecell{93\newline{}170} \\
\hline
\end{tabular}\par}

<자료 4> 본건 아파트단지의 층별 효용지수 \par
{\small\centering
\begin{tabular}{|c|c|c|c|c|c|c|c|c|c|c|}
\hline
\makecell{층} & \makecell{1} & \makecell{2} & \makecell{3} & \makecell{4} & \makecell{5} & \makecell{6} & \makecell{7} & \makecell{8} & \makecell{9} & \makecell{10} \\
\hline
\makecell{효용지수(\%)} & \makecell{90} & \makecell{94} & \makecell{96} & \makecell{98} & \makecell{99} & \makecell{100} & \makecell{100} & \makecell{100} & \makecell{100} & \makecell{100} \\
\hline
\makecell{층} & \makecell{11} & \makecell{12} & \makecell{13} & \makecell{14} & \makecell{15} & \makecell{16} & \makecell{17} & \makecell{18} & \makecell{-} & \makecell{-} \\
\hline
\makecell{효용지수(\%)} & \makecell{100} & \makecell{100} & \makecell{100} & \makecell{100} & \makecell{100} & \makecell{100} & \makecell{100} & \makecell{98} & \makecell{-} & \makecell{-} \\
\hline
\end{tabular}\par}

<자료 5> 본건 아파트단지의 위치별 효용지수 \par
{\small\centering
\begin{tabular}{|c|c|c|c|c|}
\hline
\makecell{위치} & \makecell{1호} & \makecell{2호} & \makecell{3호} & \makecell{4호} \\
\hline
\makecell{효용지수(\%)} & \makecell{98} & \makecell{100} & \makecell{98} & \makecell{96} \\
\hline
\end{tabular}\par}

<자료 6> 본건 아파트단지의 면적 타입별 효용지수 \par
{\small\centering
\begin{tabular}{|c|c|c|}
\hline
\makecell{면적(m\textsuperscript{2})} & \makecell{85} & \makecell{110} \\
\hline
\makecell{효용지수(\%)} & \makecell{100} & \makecell{104} \\
\hline
\end{tabular}\par}

<자료 7> 본건 아파트 단지내 거래사례 자료 \par
1) 창면적 감소가 없는 사례 \par
- 동, 호수 : 107동 503호 \par
- 거래시점 : 2011. 06. 25 \par
- 거래가격 : 322,000,000원 \par
2) 창면적 감소가 있는 사례 \par
- 동, 호수 : 101동 401호 \par
- 거래시점 : 2011. 03. 12 \par
- 거래가격 : 305,000,000원 \par
<자료 8> 인근지역의 아파트가격 변동지수 \par
{\small\centering
\begin{tabular}{|c|c|c|c|c|c|c|}
\hline
\makecell{2011.01.01} & \makecell{2011.02.01} & \makecell{2011.03.01} & \makecell{2011.04.01} & \makecell{2011.05.01} & \makecell{2011.06.01} & \makecell{2011.07.01} \\
\hline
\makecell{113} & \makecell{114} & \makecell{116} & \makecell{116.8} & \makecell{117.8} & \makecell{119} & \makecell{120} \\
\hline
\end{tabular}\par}

< 자료 9> 기타 평가조건 \par
1) 각 호별 정상가격 산정에 있어서 개별소유자가 개별투자한 내부마감재, 구조변경, 추가설비 및 관리상태의 차이 등의 개별적 사항은 고려하지 않는다. \par
2) 본건 아파트 단지 내 각 동별 효용격차는 없는 것으로 가정한다. \par
\end{longquestion}
